% !TeX root = ../../金融经济学_王江_讲义及习题.tex
% 02_基本框架 课后习题（由章节文件拆分）

\begin{problem}
\problemno{2.1} $U(c)$ 和 $V(c)$ 是两个效用函数，$c\in\mathrm{R}_{+}^{n}$，且 $V(x)=f(U(x))$，其中 $f(\cdot)$ 是一正单调函数。证明这两个效用函数表示了相同的偏好。
\end{problem}
\begin{solution}
定义 2.4：效用函数表示偏好 $\succeq$，是指对任意两个消费计划，$U(c_{1})\ge U(c_{2})$ 当且仅当 $c_{1}\succeq c_{2}$。两个效用函数表示相同偏好，就是它们给出完全一样的排序。正单调变换不改变不等号的方向，所以 $V=f(U)$ 与 $U$ 排序相同。

假设 $U(c)$ 表示的偏好关系为 $\succeq$，那么 $\forall\, c_{1},\,c_{2}\in\mathrm{R}_{+}^{N}$ 有
\[
    U(c_{1})\ge U(c_{2})\Leftrightarrow c_{1}\succeq c_{2}
\]
而 $f(\cdot)$ 是正单调函数，因而
\[
    V(c_{1})=f(U(c_{1}))\ge f(U(c_{2}))=V(c_{2})\Leftrightarrow U(c_{1})\ge U(c_{2})
\]
因此 $V(c_{1})\ge V(c_{2})\Leftrightarrow c_{1}\succeq c_{2}$，即 $V(c)$ 表示的偏好也是 $\succeq$。
\end{solution}

\begin{problem}
\problemno{2.2*} 在 1 期，经济有两个可能状态 $a$ 和 $b$，它们的发生概率相等：
\begin{center}
\WJEventTree{a,b}
\end{center}
考虑定义在消费计划 $c=\begin{bmatrix}c_{0} \\ c_{1a} \\ c_{1b}\end{bmatrix}$ 上的效用函数：
\[
    \begin{aligned}
        U(c) &= \log c_{0}+\tfrac12(\log c_{1a}+\log c_{1b})\\
        U(c) &= \tfrac{1}{1-\gamma}c_{0}^{1-\gamma}+\tfrac12\Bigl(\tfrac{1}{1-\gamma}c_{1a}^{1-\gamma}+\tfrac{1}{1-\gamma}c_{1b}^{1-\gamma}\Bigr)\\
        U(c) &= -e^{-ac_{0}}-\tfrac12\bigl(e^{-ac_{1a}}+e^{-ac_{1b}}\bigr)
    \end{aligned}
\]
证明它们满足：不满足性、连续性和凸性。
\end{problem}
\begin{solution}
三个效用函数都具有形式 $U(c)=u(c_{0})+\tfrac12\bigl(u(c_{1a})+u(c_{1b})\bigr)$。对数与幂效用要求底数为正，指数效用 $u(x)=-e^{-ax}$ 在 $a>0$ 时于整个实数轴上有定义。下面统一在严格正消费集
\[
    C=\bigl\{c\in\mathrm{R}^{3}:c_{0}>0,\;c_{1a}>0,\;c_{1b}>0\bigr\}
\]

上证明；并设幂效用中 $\gamma>0$、$\gamma\neq 1$，指数效用中 $a>0$。对应的期效用为
\[
    u(x)=\log x,\qquad
    u(x)=\frac{x^{1-\gamma}}{1-\gamma},\qquad
    u(x)=-e^{-ax},\qquad x>0.
\]

三者都满足 $u'(x)>0$、$u''(x)<0$，因而严格递增且严格凹：
\[
    (\log x)'=\frac1x,\quad (\log x)''=-\frac1{x^{2}},
\]
\[
    \Bigl(\frac{x^{1-\gamma}}{1-\gamma}\Bigr)'=x^{-\gamma},\quad
    \Bigl(\frac{x^{1-\gamma}}{1-\gamma}\Bigr)''=-\gamma x^{-\gamma-1},
\]
\[
    (-e^{-ax})'=ae^{-ax},\quad (-e^{-ax})''=-a^{2}e^{-ax}.
\]

\begin{enumerate}[label=(\alph*)]
    \item 不满足性（喜多厌少）。设 $c,c'\in C$ 且 $c\ge c'$。由 $u$ 在 $(0,\infty)$ 上递增，$u(c_{i})\ge u(c_{i}')$，故 $U(c)\ge U(c')$。
    \item 连续性。要证：若 $\{c^{(n)}\}\subset C$ 且 \(\lim_{n\to\infty}c^{(n)}=c\in C\)，则
    \[
        \lim_{n\to\infty}U\bigl(c^{(n)}\bigr)=U(c).
    \]

    按分量有 \(\lim_{n\to\infty}c_{i}^{(n)}=c_{i}>0\)。\(u\) 在 \((0,\infty)\) 上连续，因而在点 \(c_{i}\) 连续。由连续的序列定义，极限可以代入 \(u\)：
    \[
        \lim_{n\to\infty}u\bigl(c_{i}^{(n)}\bigr)=u\Bigl(\lim_{n\to\infty}c_{i}^{(n)}\Bigr)=u(c_{i}).
    \]
    $U$ 只是三项 $u$ 的和，极限穿过和式：
    \[
        \begin{aligned}
            \lim_{n\to\infty}U\bigl(c^{(n)}\bigr)
            &=\lim_{n\to\infty}u\bigl(c_{0}^{(n)}\bigr)
            +\tfrac12\lim_{n\to\infty}u\bigl(c_{1a}^{(n)}\bigr)
            +\tfrac12\lim_{n\to\infty}u\bigl(c_{1b}^{(n)}\bigr)\\
            &=u(c_{0})+\tfrac12\bigl(u(c_{1a})+u(c_{1b})\bigr)
            =U(c).
        \end{aligned}
    \]
    \item 凸性。$u$ 严格凹，故对 $\alpha\in(0,1)$ 与 $c,c'\in C$，
    \[
        u\bigl(\alpha c_{i}+(1-\alpha)c_{i}'\bigr)\ge\alpha u(c_{i})+(1-\alpha)u(c_{i}'),
    \]
    且凸组合仍在 $C$ 中。因此
    \[
        U\bigl(\alpha c+(1-\alpha)c'\bigr)\ge\alpha U(c)+(1-\alpha)U(c').
    \]

    若 $U(c)>U(c')$，则右端严格大于 $U(c')$，即 $U\bigl(\alpha c+(1-\alpha)c'\bigr)>U(c')$。
\end{enumerate}
\end{solution}

\begin{problem}
\problemno{2.3} $U(c)=c-\tfrac12 a c^{2}$ 是一可能的效用函数，其中 $c\in\mathrm{R}_{+}$，$a$ 是非负的系数。$U(c)$ 具有不满足性吗？如果不是，那么 $a$ 取什么值和/或 $c$ 在什么范围内时 $U(c)$ 具有不满足性？
\end{problem}
\begin{solution}
不一定，比如当 $a=1$ 时，$U(\tfrac12)=\tfrac38<U(1)=\tfrac12>U(3)=-1.5$。$U(c)$ 不具有不满足性。

当 $a=0$ 时，$U(c)=c$ 具有不满足性；

当 $a>0$ 时，当 $c\in[0,\tfrac{1}{a}]$ 时 $U(c)$ 具有不满足性。

\begin{center}
\footnotesize
\begin{tabular}{@{}ccc@{}}
\begin{tikzpicture}[baseline=(current bounding box.center), x=1cm, y=1cm]
    \path (-0.32,-1.55) rectangle (2.55,1.62);
    \draw[-{Stealth[length=1.5mm]}, thin] (0,0) -- (2.28,0) node[below left, inner sep=1pt] {$c$};
    \draw[-{Stealth[length=1.5mm]}, thin] (0,0) -- (0,1.48) node[above left, inner sep=1pt] {$U$};
    \draw[blue!65!black, thick, domain=0:2] plot (\x,{0.65*\x});
\end{tikzpicture}
&
\begin{tikzpicture}[baseline=(current bounding box.center), x=1cm, y=1cm]
    \path (-0.32,-1.55) rectangle (2.55,1.62);
    \draw[-{Stealth[length=1.5mm]}, thin] (0,0) -- (2.28,0) node[below left, inner sep=1pt] {$c$};
    \draw[-{Stealth[length=1.5mm]}, thin] (0,0) -- (0,1.48) node[above left, inner sep=1pt] {$U$};
    \draw[blue!65!black, thick, domain=0:1, samples=40, smooth] plot ({2*\x},{2.2*(\x-0.5*\x*\x)});
    \draw[dashed, black!45] (2,0) -- (2,1.1);
    \node[below, inner sep=1pt] at (2,0) {$1/a$};
\end{tikzpicture}
&
\begin{tikzpicture}[baseline=(current bounding box.center), x=1cm, y=1cm]
    \path (-0.32,-1.55) rectangle (2.55,1.62);
    \draw[-{Stealth[length=1.5mm]}, thin] (0,0) -- (2.35,0) node[above right, inner sep=1pt] {$c$};
    \draw[-{Stealth[length=1.5mm]}, thin] (0,-1.4) -- (0,0.85) node[above left, inner sep=1pt] {$U$};
    \draw[blue!65!black, thick, domain=0:3.05, samples=50, smooth] plot ({0.68*\x},{0.85*(\x-0.5*\x*\x)});
    \draw[dashed, black!45] (0.34,0) -- (0.34,0.319);
    \draw[dashed, black!45] (0.68,0) -- (0.68,0.425);
    \draw[dashed, black!45] (2.04,0) -- (2.04,-1.275);
    \fill (0.34,0.319) circle (1pt);
    \fill (0.68,0.425) circle (1pt);
    \fill (2.04,-1.275) circle (1pt);
    \draw (0.34,0.045) -- (0.34,-0.045) node[below, inner sep=1.5pt] {$\tfrac12$};
    \draw (0.68,0.045) -- (0.68,-0.045) node[below, inner sep=1.5pt] {$1$};
    \draw (2.04,0.045) -- (2.04,-0.045) node[below, inner sep=1.5pt] {$3$};
\end{tikzpicture}
\\[0.15em]
$a=0$，$U(c)=c$ & $a>0$，$c\in[0,1/a]$ & $a=1$，$U(\tfrac12)<U(1)>U(3)$
\end{tabular}
\end{center}
\end{solution}

\begin{problem}
\problemno{2.4} 考虑一个经济，它在 1 期有三个可能状态：$a$，$b$ 和 $c$：
\begin{center}
\WJEventTree{a,b,c}
\end{center}
证券市场包括证券 1 和 2，它们具有如下的支付向量：$X_{1}=\begin{bmatrix}1 \\ 1 \\ 1\end{bmatrix}$ 以及 $X_{2}=\begin{bmatrix}1 \\ 2 \\ 3\end{bmatrix}$。它们的价格分别为 $S_{1}$ 和 $S_{2}$。
\begin{enumerate}[label=(\alph*)]
    \item 描述这个经济的支付空间。
    \item 写出这个经济的市场结构矩阵 $X$。
    \item 考虑含有 $\theta_{1}$ 单位的证券 1 和 $\theta_{2}$ 单位的证券 2 的组合。写出这个组合的支付向量。这个组合的价格是多少？
    \item 假设这个市场中总共有 $K$ 个参与者。每个参与者的禀赋是 1 单位的证券 1 和 2 单位的证券 2。这时的市场组合是什么？市场组合的支付向量是什么？市场组合的总价值是多少？
    \item 写出市场化支付的集合。
    \item 如果市场不允许卖空。市场化支付的集合是什么？
    \item 现在引入新的证券 3，它的支付向量为 $X_{3}=\begin{bmatrix}0 \\ 0 \\ 1\end{bmatrix}$。写出新的市场结构矩阵。在这个市场结构下，市场化支付集合是什么？
\end{enumerate}
\end{problem}
\begin{solution}
三个状态 $a,b,c$ 对应 $\Omega=3$。记
\[
    X_1=\begin{bmatrix}1 \\ 1 \\ 1\end{bmatrix},\quad
    X_2=\begin{bmatrix}1 \\ 2 \\ 3\end{bmatrix},\quad
    \theta=\begin{bmatrix}\theta_1 \\ \theta_2\end{bmatrix}.
\]

组合在 1 期的支付向量为 $X_1\theta_1+X_2\theta_2$，也可写成矩阵乘积 $X\theta$，其中 $X=[X_1,\,X_2]$。

\begin{enumerate}[label=(\alph*)]
    \item \textbf{支付空间}：所有可能在三个状态下实现的支付向量的集合，即 $\mathrm{R}^{\Omega}=\mathrm{R}^{3}$（每个分量对应状态 $a,b,c$ 下的支付）。

    \item \textbf{市场结构矩阵}：两列分别为两只证券的支付向量，
    \[
        X=\begin{bmatrix}1&1\\1&2\\1&3\end{bmatrix}=[X_1,\,X_2].
    \]

    \item \textbf{组合支付与价格：}持有 $(\theta_1,\theta_2)$ 时，
    \[
        X\theta=\theta_1 X_1+\theta_2 X_2=
        \begin{bmatrix}\theta_1+\theta_2 \\ \theta_1+2\theta_2 \\ \theta_1+3\theta_2\end{bmatrix}.
    \]

    0 期购买成本（组合价格）为 $\theta_1 S_1+\theta_2 S_2$。

    \item \textbf{市场组合：}$K$ 个参与者每人禀赋为 $1$ 单位证券 1、$2$ 单位证券 2，加总得市场组合
    \[
        \theta_M=\begin{bmatrix}K \\ 2K\end{bmatrix}.
    \]

    其支付向量为
    \[
        X\theta_M=K X_1+2K X_2=
        \begin{bmatrix}3K \\ 5K \\ 7K\end{bmatrix}.
    \]

    总价值为 $K S_1+2K S_2$。

    \item \textbf{市场化支付（允许卖空）：}能由证券组合复制的 1 期支付集合为
    \[
        M=\bigl\{Y\in\mathrm{R}^{3}: Y=X\theta=[X_1,\,X_2]\begin{bmatrix}\theta_1 \\ \theta_2\end{bmatrix}=\theta_1 X_1+\theta_2 X_2,\;\theta_1,\theta_2\in\mathrm{R}\bigr\},
    \]

    即 $X_1,X_2$ 张成的二维子空间（只有两只证券，维度 $\le 2$）。

    \item \textbf{禁止卖空：}仅允许 $\theta_1,\theta_2\ge 0$，为 $X_1,X_2$ 的\textbf{非负}线性组合：
    \[
        M_{+}=\bigl\{Y\in\mathrm{R}^{3}: Y=X\theta=[X_1,\,X_2]\begin{bmatrix}\theta_1 \\ \theta_2\end{bmatrix}=\theta_1 X_1+\theta_2 X_2,\;\theta_1,\theta_2\in\mathrm{R}_{+}\bigr\}.
    \]

    \item \textbf{加入证券 3：}其支付向量为 $X_3=\begin{bmatrix}0 \\ 0 \\ 1\end{bmatrix}$，仅在状态 $c$ 支付 $1$。两只证券时的市场结构是 $X=[X_1,\,X_2]$。加入证券 3 后的市场结构另记为 $X'$，以免与 $X$ 相混：
    \[
        X'=\begin{bmatrix}1&1&0\\1&2&0\\1&3&1\end{bmatrix}=[X_1,\,X_2,\,X_3].
    \]

    三列线性无关，可表示任意 $Y\in\mathrm{R}^{3}$，故（允许卖空时）
    \[
        M=\bigl\{Y\in\mathrm{R}^{3}: Y=\theta_1 X_1+\theta_2 X_2+\theta_3 X_3,\;\theta_1,\theta_2,\theta_3\in\mathrm{R}\bigr\}=\mathrm{R}^{3}.
    \]
\end{enumerate}
\end{solution}

\begin{problem}
\problemno{2.5} 在练习 2.4 中定义的只存在证券 1 和 2 的经济中。考虑一个禀赋为 $\theta_{1}$ 单位的证券 1 和 $\theta_{2}$ 单位的证券 2 的参与者。写出他的预算集。
\end{problem}
\begin{solution}
无摩擦，允许卖空。预算集是花费不超过禀赋市值的那些 1 期支付：
\[
    B_{12}=\Bigl\{C\in\mathrm{R}_{+}^{3}:
    C=\alpha_{1}X_{1}+\alpha_{2}X_{2},\;
    \alpha_{1}S_{1}+\alpha_{2}S_{2}\le\theta_{1}S_{1}+\theta_{2}S_{2},\;
    \alpha_{1},\alpha_{2}\in\mathrm{R}\Bigr\}.
\]
\end{solution}

\begin{problem}
\problemno{2.6} 在上面的练习中引入练习 2.4 中定义的证券 3，它的价格为 $S_{3}$。这时，参与者的预算集是什么（他在证券 3 上的禀赋为 0）？证明由证券 1、2、3 构成的预算集包含仅由证券 1、2 构成的预算集。
\end{problem}
\begin{solution}
证券 3 上的禀赋为 $0$，可交易财富仍是 $\theta_{1}S_{1}+\theta_{2}S_{2}$。预算集为
\[
    B_{123}=\Bigl\{C\in\mathrm{R}_{+}^{3}:
    C=\alpha_{1}X_{1}+\alpha_{2}X_{2}+\alpha_{3}X_{3},\;
    \alpha_{1}S_{1}+\alpha_{2}S_{2}+\alpha_{3}S_{3}\le\theta_{1}S_{1}+\theta_{2}S_{2},\;
    \alpha_{1},\alpha_{2},\alpha_{3}\in\mathrm{R}\Bigr\}.
\]

练习 2.5 的预算集是
\[
    B_{12}=\Bigl\{C\in\mathrm{R}_{+}^{3}:
    C=\alpha_{1}X_{1}+\alpha_{2}X_{2},\;
    \alpha_{1}S_{1}+\alpha_{2}S_{2}\le\theta_{1}S_{1}+\theta_{2}S_{2},\;
    \alpha_{1},\alpha_{2}\in\mathrm{R}\Bigr\}.
\]

设 $C\in B_{12}$，对应头寸为 $\alpha_{1},\alpha_{2}$。取 $\alpha_{3}=0$，则
\[
    C=\alpha_{1}X_{1}+\alpha_{2}X_{2}+\alpha_{3}X_{3},
\]

并且
\[
    \alpha_{1}S_{1}+\alpha_{2}S_{2}+\alpha_{3}S_{3}=\alpha_{1}S_{1}+\alpha_{2}S_{2}\le\theta_{1}S_{1}+\theta_{2}S_{2}.
\]

故 $C\in B_{123}$，即 $B_{12}\subseteq B_{123}$。
\end{solution}

\begin{problem}
\problemno{2.7*} 考虑一个在 1 期只有一个可能状态的经济。（在这种情况下不存在不确定性。）参与者 1 的 0 期禀赋为 100 而 1 期禀赋为 1，即他的禀赋向量为 $\begin{bmatrix}100 \\ 1\end{bmatrix}$。他的偏好可以表示成如下形式：
\[
    U(c_{0},c_{1})=\log c_{0}+\rho\log c_{1}.
\]

系数 $\rho$ 为反映参与者在当前消费和未来消费之间相对偏好的参数。有一只证券，它的 0 期价格为 1、1 期支付为 $1+r_{F}$。这里，$r_{F}$ 是利率。
\begin{enumerate}[label=(\alph*)]
    \item 如果这个参与者不能在市场上进行交易，那么他的消费计划以及相应的效用 $U^{a}$ 是什么？
    \item 现在假设他可以在市场上进行交易。
    \begin{itemize}
        \item 他的预算集是什么？以当前消费为单位，他的总财富 $w$ 是多少？
        \item 写出参与者的优化问题。令 $c_{0}$ 为参与者的当前（即 0 期）最优消费，$s$ 为最优储蓄以及 $U^{b}$ 为在最优策略下得到的效用。求解他的最优消费/储蓄选择以及相应的效用。把 $U^{b}$ 表示成财富 $w$、利率 $r_{F}$ 和偏好系数 $\rho$ 的函数。
        \item 讨论参与者的最优选择如何依赖于利率 $r_{F}$ 和偏好系数 $\rho$。给出解释。
    \end{itemize}
    \item 证明 $U^{b}\ge U^{a}$。
    \item 令 $g$ 为参与者由于能够在证券市场上交易而获得的益处。它的定义为
    \[
        U^{b}(w-g)=U^{a}.
    \]

    计算 $g$。讨论 $g$ 如何依赖于 $\rho$？$g$ 如何依赖于 $r_{F}$？给出解释。
\end{enumerate}
\end{problem}
\begin{solution}
禀赋 $e_0=100$、$e_1=1$。无风险证券：0 期价格 $1$，1 期支付 $1+r_F$。记 $S$ 为 0 期\textbf{储蓄}（买入该证券的金额；$S<0$ 表示借贷）。

\begin{enumerate}[label=(\alph*)]
    \item \textbf{不能交易（自给自足）：}只能消费禀赋，
    \[
        c_0=100,\qquad c_1=1,\qquad
        U^{a}=\log(100)+\rho\log(1)=\log 100.
    \]

    \item \textbf{可以交易。}
    \begin{itemize}
        \item \textbf{预算集与财富 $w$。}买卖证券等价于在两期之间转移消费：0 期少消费 $S$ 用于储蓄，1 期增加 $S(1+r_F)$。故
        \[
            c_0=100-S,\qquad c_1=1+S(1+r_F),\qquad S\in\mathrm{R}.
        \]

        以 0 期消费为计价单位，把 1 期禀赋折现到 0 期（用市场利率 $r_F$）：
        \[
            w=100+\frac{1}{1+r_F}.
        \]

        这是参与者按市价出售全部禀赋后可在 0 期支配的\textbf{总财富}（现值）。任意满足上述两式的 $(c_0,c_1)$ 都在预算集内。

        \item \textbf{优化问题。}在预算集上最大化对数效用：
        \[
            \max_{S\in\mathrm{R}}\ \log(100-S)+\rho\log\bigl[1+S(1+r_F)\bigr].
        \]

        一阶条件（边际效用替代率 = 跨期相对价格）：
        \[
            \underbrace{\frac{1}{100-S}}_{\text{0 期边际效用}}
            =\rho\underbrace{\frac{1+r_F}{1+S(1+r_F)}}_{\text{1 期边际效用}\times\text{价格比}}.
        \]

        整理得 $\bigl[1+S(1+r_F)\bigr]=\rho(1+r_F)(100-S)$，即
        \[
            1+S(1+r_F)(1+\rho)=100\rho(1+r_F).
        \]
        
        其解记为最优储蓄
        \[
            S^{*}=\frac{100\rho(1+r_F)-1}{(1+\rho)(1+r_F)}.
        \]

        预算约束是前面的两式 \(c_0=100-S\)、\(c_1=1+S(1+r_F)\)，其中 \(c_0,c_1,S\) 仍是可行选择，不是最优值。把 \(S^{*}\) 代回去，得到最优消费 \(c_0^{*},c_1^{*}\)：
        \[
            \begin{aligned}
                c_0^{*}&=100-S^{*}
                =100-\frac{100\rho(1+r_F)-1}{(1+\rho)(1+r_F)}
                =\frac{100(1+r_F)+1}{(1+\rho)(1+r_F)},\\[0.4em]
                c_1^{*}&=1+S^{*}(1+r_F)
                =1+(1+r_F)\cdot\frac{100\rho(1+r_F)-1}{(1+\rho)(1+r_F)}
                =\frac{\rho\bigl[100(1+r_F)+1\bigr]}{1+\rho}.
            \end{aligned}
        \]

        财富 \(w=100+\dfrac{1}{1+r_F}=\dfrac{100(1+r_F)+1}{1+r_F}\) 的分子正是上式的 \(100(1+r_F)+1\)，故
        \[
            c_0^{*}=\frac{w}{1+\rho},\qquad
            c_1^{*}=\frac{\rho(1+r_F)}{1+\rho}\,w.
        \]

        两式相除，零期与一期的最优消费只差偏好权重和总收益：
        \[
            c_1^{*}=\rho(1+r_F)\,c_0^{*}.
        \]

        因此最优现期消费占财富 $\dfrac{1}{1+\rho}$，这个份额不随 \(r_F\) 改变；最优 1 期消费占 \(\dfrac{\rho(1+r_F)}{1+\rho}\)，是在该份额上再乘以 \(\rho(1+r_F)\)。这是对数效用下的固定支出份额。最优效用
        \[
            U^{b}=\log\Bigl(\frac{w}{1+\rho}\Bigr)+\rho\log\Bigl(\frac{\rho(1+r_F)}{1+\rho}\,w\Bigr).
        \]

        \item \textbf{比较静态。}上面已经求出最优消费量：
        \[
            c_0^{*}=\frac{100(1+r_F)+1}{(1+\rho)(1+r_F)}
            =\frac{1}{1+\rho}\Bigl(100+\frac{1}{1+r_F}\Bigr),\qquad
            c_1^{*}=\frac{\rho\bigl[100(1+r_F)+1\bigr]}{1+\rho}.
        \]
        \(r_F\) 只出现在 \(c_0^{*}\) 的 \(1/(1+r_F)\) 里，故 \(r_F\) 上升时 \(c_0^{*}\) 下降；\(\rho\) 在分母，\(\rho\) 上升时 \(c_0^{*}\) 也下降。
        
        \(c_1^{*}\) 的分子含 \(\rho(1+r_F)\)，分母是 \(1+\rho\)，两项都使 \(c_1^{*}\) 随 \(r_F\)、\(\rho\) 上升。
        
        直觉：\(r_F\) 升高提高储蓄回报，\(\rho\) 升高加重未来消费权重，都会促使少消费于 0 期、多留于 1 期。
    \end{itemize}

    \item \textbf{$U^{b}\ge U^{a}$。}不能交易时只能消费禀赋，\((c_0,c_1)=(100,1)\)，效用为 \(U^{a}=\log 100\)。
    
    能交易时储蓄 \(S\) 可任取，其中 \(S=0\) 仍可行，对应的消费还是 \((100,1)\)，效用仍是 \(\log 100=U^{a}\)。\(U^{b}\) 是所有可行储蓄里的最大效用，而 \(S=0\) 只是其中一个可行点，最大值不会低于这一点：
    \[
        U^{b}=\max_{S\in\mathrm{R}}\,U\bigl(100-S,\,1+S(1+r_F)\bigr)
        \ge U(100,1)=U^{a}.
    \]
    等号当且仅当最优储蓄恰好是 \(S^{*}=0\)，即 \(100\rho(1+r_F)=1\)。

    \item \textbf{交易益处 $g$。}定义 $U^{b}(w-g)=U^{a}$：$w-g$ 是与交易后最优效用相同、但用一条 0 期财富数字表示的\textbf{确定性等价财富}；$g=w-(w-g)$ 即为因能交易而多出来的等价 0 期资源。

    将 $U^{a}=\log 100$ 与上式 $U^{b}$ 联立，记 $w-g=\hat w$，则
    \[
        (1+\rho)\log\hat w+\rho\log\bigl(\rho(1+r_F)\bigr)-(1+\rho)\log(1+\rho)=\log 100.
    \]

    取指数并整理：
    \[
        \hat w^{1+\rho}\,\frac{\bigl[\rho(1+r_F)\bigr]^{\rho}}{(1+\rho)^{1+\rho}}=100
        \quad\Rightarrow\quad
        \hat w=\rho^{-\frac{\rho}{1+\rho}}(1+r_F)^{-\frac{\rho}{1+\rho}}
        100^{\frac{1}{1+\rho}}(1+\rho).
    \]

    故
    \[
        g=w-\hat w
        =w-\rho^{-\frac{\rho}{1+\rho}}(1+r_F)^{-\frac{\rho}{1+\rho}}
        100^{\frac{1}{1+\rho}}(1+\rho).
    \]
    $g$ 随 $\rho$、$r_F$ 增加而增加：二者都提高把消费从 0 期转到 1 期」的动机，能在市场上做跨期配置带来的等价收益更大。
\end{enumerate}
\end{solution}

